\documentclass[11pt]{article}
\usepackage[letterpaper,margin=1in]{geometry}
\usepackage{amsmath,amssymb,booktabs,array}
\usepackage[T1]{fontenc}
\usepackage{lmodern}
\usepackage{xcolor}
\usepackage{listings}
\usepackage[hidelinks]{hyperref}
\usepackage{fancyhdr}
\usepackage{microtype}

\definecolor{softgray}{gray}{0.94}
\lstset{
  language=Python,
  basicstyle=\ttfamily\footnotesize,
  keywordstyle=\bfseries,
  commentstyle=\itshape,
  stringstyle=\ttfamily,
  backgroundcolor=\color{softgray},
  frame=single,
  rulecolor=\color{black},
  numbers=left,
  numberstyle=\scriptsize,
  breaklines=true,
  showstringspaces=false,
  columns=fullflexible
}

\pagestyle{fancy}
\fancyhf{}
\lhead{Homework 1}
\rhead{Matrix Algebra}
\cfoot{\thepage}
\setlength{\headheight}{14pt}

\title{\textbf{Homework 1: Matrix Algebra}\\[4pt]\large Manual Solutions and NumPy Verification}
\author{Xinchen Leng}
\date{September 29, 2026}

\begin{document}
\maketitle

\begin{center}
\small
\textbf{Code repository:}\quad
\href{https://github.com/Orlandoxinchen/homework1-linear-algebra}
{github.com/Orlandoxinchen/homework1-linear-algebra}
\end{center}

\noindent\rule{\textwidth}{0.6pt}

\section*{Part A. Conceptual / Manual Problems}

\subsection*{1. Basic Operations}
Let
\[
A=\begin{bmatrix}1&2\\3&4\end{bmatrix},\qquad
B=\begin{bmatrix}2&0\\1&2\end{bmatrix}.
\]

\paragraph{(a) Addition and subtraction.}
Elementwise computation gives
\[
A+B=\begin{bmatrix}1+2&2+0\\3+1&4+2\end{bmatrix}
=\boxed{\begin{bmatrix}3&2\\4&6\end{bmatrix}},
\]
\[
A-B=\begin{bmatrix}1-2&2-0\\3-1&4-2\end{bmatrix}
=\boxed{\begin{bmatrix}-1&2\\2&2\end{bmatrix}}.
\]

\paragraph{(b) Scalar multiplication.}
\[
2A=\boxed{\begin{bmatrix}2&4\\6&8\end{bmatrix}}.
\]

\paragraph{(c) Matrix products.}
\[
AB=
\begin{bmatrix}
1(2)+2(1)&1(0)+2(2)\\
3(2)+4(1)&3(0)+4(2)
\end{bmatrix}
=\boxed{\begin{bmatrix}4&4\\10&8\end{bmatrix}},
\]
\[
BA=
\begin{bmatrix}
2(1)+0(3)&2(2)+0(4)\\
1(1)+2(3)&1(2)+2(4)
\end{bmatrix}
=\boxed{\begin{bmatrix}2&4\\7&10\end{bmatrix}}.
\]
Thus, $AB\ne BA$. Matrix multiplication is generally not commutative; changing the order changes which rows are combined with which columns.

\subsection*{2. Properties of Special Matrices}
The following are four different examples:
\[
\text{Diagonal: }\begin{bmatrix}3&0&0\\0&-2&0\\0&0&5\end{bmatrix},\qquad
\text{Identity: }\begin{bmatrix}1&0\\0&1\end{bmatrix},
\]
\[
\text{Symmetric: }\begin{bmatrix}2&-1\\-1&4\end{bmatrix},\qquad
\text{Idempotent: }\begin{bmatrix}1&0\\0&0\end{bmatrix}.
\]
The symmetric example equals its transpose. The idempotent example satisfies $P^2=P$.

\subsection*{3. System of Equations}
The system
\[
2x+y=5,\qquad 3x-y=4
\]
has matrix form
\[
\underbrace{\begin{bmatrix}2&1\\3&-1\end{bmatrix}}_{A}
\underbrace{\begin{bmatrix}x\\y\end{bmatrix}}_{\mathbf{x}}
=
\underbrace{\begin{bmatrix}5\\4\end{bmatrix}}_{\mathbf{b}}.
\]
Since $\det(A)=2(-1)-1(3)=-5\ne0$,
\[
A^{-1}=\frac{1}{-5}\begin{bmatrix}-1&-1\\-3&2\end{bmatrix}
=\begin{bmatrix}\tfrac15&\tfrac15\\\tfrac35&-\tfrac25\end{bmatrix}.
\]
Therefore,
\[
\begin{bmatrix}x\\y\end{bmatrix}
=A^{-1}\mathbf b
=\begin{bmatrix}\tfrac15&\tfrac15\\\tfrac35&-\tfrac25\end{bmatrix}
\begin{bmatrix}5\\4\end{bmatrix}
=\boxed{\begin{bmatrix}\tfrac95\\\tfrac75\end{bmatrix}}.
\]
Hence $(x,y)=(1.8,1.4)$.

\subsection*{4. Determinant and Rank}
For
\[
C=\begin{bmatrix}1&2&3\\1&1&1\\3&5&7\end{bmatrix},
\]
cofactor expansion along the first row gives
\begin{align*}
\det(C)
&=1\begin{vmatrix}1&1\\5&7\end{vmatrix}
-2\begin{vmatrix}1&1\\3&7\end{vmatrix}
+3\begin{vmatrix}1&1\\3&5\end{vmatrix}\\
&=(7-5)-2(7-3)+3(5-3)=2-8+6=\boxed{0}.
\end{align*}
Also, $R_3=2R_1+R_2$, so the three rows are linearly dependent. The first two rows are not scalar multiples, so they are independent. Therefore,
\[
\boxed{\operatorname{rank}(C)=2}.
\]
The rank tells us that only two rows (and two columns) provide independent information. The matrix maps vectors into a two-dimensional subspace and is singular, consistent with its zero determinant.

\subsection*{5. Matrix Inverse}
Let
\[
D=\begin{bmatrix}2&1&0\\1&1&1\\0&1&1\end{bmatrix}.
\]
Its determinant is $-1$, so it is invertible. Row reduction of $[D\mid I]$ (or the adjugate formula) gives
\[
\boxed{D^{-1}=\begin{bmatrix}0&1&-1\\1&-2&2\\-1&2&-1\end{bmatrix}}.
\]
Indeed,
\[
DD^{-1}=\begin{bmatrix}1&0&0\\0&1&0\\0&0&1\end{bmatrix}.
\]

\section*{Part B. Python Verification}
The required checks were performed with NumPy. The scripts are separated by question so each component can be run and reviewed independently.

\subsection*{Question 1: \texttt{q1\_basic\_operations.py}}
\begin{lstlisting}
import numpy as np

A = np.array([[1, 2], [3, 4]], dtype=int)
B = np.array([[2, 0], [1, 2]], dtype=int)

print("A + B =\n", A + B)
print("\nA - B =\n", A - B)
print("\n2A =\n", 2 * A)
print("\nAB =\n", A @ B)
print("\nBA =\n", B @ A)
print("\nAB equals BA:", np.array_equal(A @ B, B @ A))
\end{lstlisting}

\subsection*{Question 3: \texttt{q3\_system.py}}
\begin{lstlisting}
import numpy as np

A = np.array([[2.0, 1.0], [3.0, -1.0]])
b = np.array([5.0, 4.0])

A_inverse = np.linalg.inv(A)
solution = A_inverse @ b

print("A inverse =\n", A_inverse)
print("\n[x, y] =", solution)
print("\nCheck A @ solution =", A @ solution)
\end{lstlisting}

\subsection*{Question 4: \texttt{q4\_determinant\_rank.py}}
\begin{lstlisting}
import numpy as np

C = np.array([[1.0, 2.0, 3.0],
              [1.0, 1.0, 1.0],
              [3.0, 5.0, 7.0]])

determinant = np.linalg.det(C)
rank = np.linalg.matrix_rank(C)

print("det(C) =", round(determinant))
print("rank(C) =", rank)
print("Row 3 equals 2*Row 1 + Row 2:",
      np.allclose(C[2], 2 * C[0] + C[1]))
\end{lstlisting}

\subsection*{Question 5: \texttt{q5\_inverse.py}}
\begin{lstlisting}
import numpy as np

D = np.array([[2.0, 1.0, 0.0],
              [1.0, 1.0, 1.0],
              [0.0, 1.0, 1.0]])

D_inverse = np.linalg.inv(D)

print("D inverse =\n", D_inverse)
print("\nCheck D @ D_inverse =\n",
      np.round(D @ D_inverse, 10))
\end{lstlisting}

\section*{Verification Summary}
\begin{center}
\begin{tabular}{@{}lll@{}}
\toprule
Question & NumPy operation & Verified result \\
\midrule
1 & Array arithmetic and \texttt{@} & $AB\ne BA$ \\
3 & \texttt{np.linalg.inv} & $(x,y)=(1.8,1.4)$ \\
4 & \texttt{det}, \texttt{matrix\_rank} & $\det(C)=0$, rank$(C)=2$ \\
5 & \texttt{np.linalg.inv} & $DD^{-1}=I_3$ \\
\bottomrule
\end{tabular}
\end{center}

\vfill
\noindent\rule{\textwidth}{0.4pt}
\small\noindent The complete source code and this report are available at
\href{https://github.com/Orlandoxinchen/homework1-linear-algebra}
{https://github.com/Orlandoxinchen/homework1-linear-algebra}.

\end{document}
